\section{Cross-Market Mapping and Hypotheses}

\subsection{The unavoidable mapping problem}

An equity index has a natural market-capitalization weighting. Futures open interest is
not corporate ownership, and the number of outstanding contracts cannot be compared
across products without accounting for price units and multipliers. We therefore use two
market-return proxies that are defensible for different reasons.

The primary proxy weights each available product return by lagged open-interest notional:
\begin{equation}
  w^{OI}_{i,d}=
  \frac{P_{i,d-1}M_iOI_{i,d-1}}
  {\sum_j P_{j,d-1}M_jOI_{j,d-1}},
  \qquad
  r^{M,OI}_{d}=\sum_i w^{OI}_{i,d}r_{i,d},
  \label{eq:oi-market}
\end{equation}
where $P$, $M$, and $OI$ denote lagged price, contract multiplier, and lagged open
interest. This proxy emphasizes products with larger outstanding notional activity. The
robustness proxy assigns equal weight to each available product return. It emphasizes the
typical product rather than the concentration of notional activity.

Neither proxy is declared universally correct. Their disagreement is part of the test.
Choosing the proxy that produces the best endpoint after observing returns would replace
cross-market replication with specification search.

\subsection{Pre-specified hypotheses}

The source mechanism implies three directional hypotheses:

\begin{enumerate}
  \item Momentum should dominate reversal in volatility states Q1--Q3.
  \item Reversal should strengthen and dominate momentum in Q4.
  \item The Q4 switch should improve momentum, reversal, and a static equal blend without
  depending on one event or one market proxy.
\end{enumerate}

The third condition is stricter than comparing the router only with reversal. A selector
can appear successful simply because it spends most months in the stronger sleeve. It
must add information relative to the best standalone sleeve and to a no-state
diversification benchmark.

\subsection{Why the static benchmark is 50/50}

Momentum and reversal are constructed with the same unit long and unit short gross
exposure. Their static blend is
\begin{equation}
  R^{50/50}_{m}=0.5R^{MOM}_{m}+0.5R^{REV}_{m}.
\end{equation}
The equal allocation is not estimated. It is the parameter-free question: does dynamic
state selection improve on holding both sleeves without forecasting the state? It is not
claimed to be the optimal ex-post combination.
